\begin{table}[htbp]
\centering
\small
\caption{Калибровка вероятностей и обмен покрытия на риск}
\label{tab:6}
\begin{tabular}{lllll}
\toprule
Величина & Macro-среднее по 18 holdout & Размах по holdout & Худший holdout & Статус величины \\
\midrule
Brier & 0.0365 & 0.0109–0.1796 & genre\_seo & зарегистрированная метрика \\
ECE, 10 бинов & 0.0571 & 0.0113–0.1973 & genre\_seo & описательная \\
MCE, 10 бинов & 0.2647 & 0.0523–0.7548 & genre\_seo & описательная \\
Риск при покрытии 100\% & 0.0468 & 0.0173–0.2014 & genre\_seo & описательная \\
Риск при покрытии 90\% & 0.0189 & 0.0000–0.1774 & genre\_seo & описательная \\
Риск при покрытии 80\% & 0.0124 & 0.0000–0.1425 & genre\_seo & описательная \\
Риск при покрытии 70\% & 0.0085 & 0.0000–0.1076 & genre\_seo & описательная \\
Риск при покрытии 50\% & 0.0059 & 0.0000–0.0854 & genre\_seo & описательная \\
\bottomrule
\end{tabular}
\begin{minipage}{\linewidth}\footnotesize\vspace{2pt}
\textit{Примечание.} Единица анализа — документ внутри группового holdout. Знаменатель: 18 групповых holdout; macro-среднее берётся по holdout, а не по документам, поэтому крупные holdout не перевешивают мелкие. Данные: серия calibration-v1, посчитана read-only по предсказаниям fairness-v2. Brier — средний квадрат отклонения вероятности от исхода; ECE и MCE — средний и максимальный разрыв между вероятностью и частотой на равночастотных бинах; риск — доля ошибочных решений среди оставленных при сокращении покрытия. Интервалы не считались. Шкала: меньше значит лучше у всех величин таблицы. Сокращения: ECE — expected calibration error; MCE — maximum calibration error; покрытие — доля документов, по которым решение принимается, остальные уходят в отказ. Пропуски: пропусков нет; порог по этой таблице не выбирается. Поправка на множественные сравнения не применялась.
\end{minipage}
\end{table}
